\section{Postfix to Infix Conversion}
\label{sec:sq-postfix-to-infix}

\problemstatement{Given a postfix expression, convert it to infix
notation (fully parenthesized, to remove any ambiguity).}

\begin{patternbox}
Converting \emph{from} postfix or prefix is structurally easier than
converting \emph{from} infix (Sections~\ref{sec:sq-infix-to-postfix}--\ref{sec:sq-infix-to-prefix}),
because postfix and prefix already encode the expression tree's structure
unambiguously through ordering alone -- there is no precedence or
associativity to resolve. The keyword shift to notice: when the
\emph{source} notation is postfix or prefix, the stack holds
\textbf{partially built sub-expressions} (strings), not raw operators, and
every operator immediately and unambiguously tells you which two
sub-expressions to combine.
\end{patternbox}

\subsection*{Understanding the problem carefully}

In postfix, by the time we encounter an operator while scanning left to
right, both of its operands have already been fully read and are sitting
on top of the stack (the most recent two pushes) -- there is never any
ambiguity about which sub-expressions an operator applies to, unlike
infix where precedence must be consulted.

\begin{ideabox}[Stack of Sub-expressions]
Scan left to right. Push every operand as a (string) sub-expression. On
an operator: pop the top two sub-expressions -- call the first popped
$\textit{right}$ and the second popped $\textit{left}$ (since
$\textit{left}$ was pushed earlier and so sits one level deeper). Combine
them as $(\textit{left}\ \textit{op}\ \textit{right})$, fully
parenthesized, and push this combined string back as a single
sub-expression.
\end{ideabox}

\subsection*{Why the most recent two pushes are always the correct operands}

Postfix notation guarantees that an operator always immediately follows
the complete sub-expressions it operates on -- this is precisely what
``postfix'' (operator \emph{after} operands) means. So whatever is on top
of the stack at the moment we read an operator must be its right operand
(read most recently), and whatever is just below that must be its left
operand, with no possibility of some other, older sub-expression
interfering, because every operator so far has already fully consumed and
replaced its own operands before we reach this point.

\subsection*{Algorithm}

\begin{enumerate}
  \item Initialize an empty stack of strings.
  \item For each token in the postfix expression:
  \begin{enumerate}
    \item Operand: push it as a one-character string.
    \item Operator $\textit{op}$: pop $\textit{right}$, pop
          $\textit{left}$; push $(\textit{left}\ \textit{op}\
          \textit{right})$.
  \end{enumerate}
  \item The stack's single remaining element is the infix result.
\end{enumerate}

\subsection*{Dry run}

\dryrun{Take $\texttt{"ab+c*"}$.}

\begin{itemize}
  \item \texttt{a}: push. Stack $=[\texttt{a}]$.
  \item \texttt{b}: push. Stack $=[\texttt{a},\texttt{b}]$.
  \item \texttt{+}: pop right$=\texttt{b}$, pop left$=\texttt{a}$. Push
        \texttt{(a+b)}. Stack $=[\texttt{(a+b)}]$.
  \item \texttt{c}: push. Stack $=[\texttt{(a+b)},\texttt{c}]$.
  \item \texttt{*}: pop right$=\texttt{c}$, pop left$=\texttt{(a+b)}$.
        Push \texttt{((a+b)*c)}.
\end{itemize}

Result: \texttt{((a+b)*c)}.

\subsection*{C++ implementation}

\begin{lstlisting}
#include <bits/stdc++.h>
using namespace std;

string postfixToInfix(string s) {
    stack<string> st;

    for (char c : s) {
        if (isalnum(c)) {
            st.push(string(1, c));
        } else {
            string right = st.top(); st.pop();
            string left = st.top(); st.pop();
            st.push("(" + left + string(1, c) + right + ")");
        }
    }
    return st.top();
}

int main() {
    cout << postfixToInfix("ab+c*") << endl;
    return 0;
}
\end{lstlisting}

\begin{complexitybox}
\textbf{Time Complexity.} Each token triggers $O(1)$ stack operations,
but string concatenation can cost up to $O(n)$ per operator in the worst
case (building progressively longer strings), giving an overall worst
case of
\[
  O(n^2)
\]
for naive string concatenation (this can be improved to $O(n)$ with more
careful string-building techniques, but $O(n^2)$ is what the straightforward
implementation above gives).
\textbf{Space Complexity.} $O(n)$ for the stack of sub-expressions.
\end{complexitybox}

\begin{takeawaybox}
When converting \emph{from} postfix or prefix, the stack holds
sub-expressions (strings), not bare operators -- every operator
encountered has an unambiguous, immediate pair of operands to combine,
with no precedence comparison ever needed. This is structurally simpler
than converting \emph{from} infix, and the same ``stack of
sub-expressions'' idea reused, with only the combination format
(parenthesized infix here, vs.\ prefix or postfix in the next few
sections) changing.
\end{takeawaybox}
